\input amstex
\documentstyle{amsppt}
\topmatter
\Monograph
\title Scrabble Pairings \endtitle
\author John J. Chew, III \endauthor
\address 
Department~of~Mathematics,
University~of~Toronto,\newline
\indent Toronto,~Ontario,
Canada~M5S~3G3
\endaddress
\email jjchew\@math.utoronto.ca\endemail
\date November 12, 2006
% (2nd ed.) 
\enddate

\keywords
pairings,
Scrabble
\endkeywords

\abstract
To follow.
\endabstract

\endtopmatter
\document
\head Introduction \endhead

$N_p$ players compete in a tournament to determine an ultimate
victor.
In each of $N_r$ rounds, a player may play a game with at most
one other player.
If players $i$ and $j$ play a game, the outcome or score
is the integer $n$ with probability $P_{ij}(n)$ and distribution 
$D_{ij}(n):=\sum_{i=-\infty}^n P_{ij}(n)$.
If $n>0$, then we say that player $i$ wins or scores a win,
and player $j$ loses or scores a loss.
If $n<0$, then player $i$ loses and player $j$ wins.
If $n=0$, then player $i$ and player $j$ both tie, scoring half a win and half
a loss.
Player $i$ therefore wins against player $j$ with probability 
$1-D_{ij}(0)$, ties with probability $P_{ij}(0)$ and loses with
probability $D_{ij}(-1)$.
If a player scores $w$ wins, $l$ losses and total score $s$ in the
first $n$ rounds, we say that the player has a record after round
$n$ of $(w,l,s)$.

Players with more wins are ranked ahead of players with fewer.
Among players with the same number of wins, those with fewer losses
are ranked ahead of those with more.
Among players with the same number of wins and losses, those with
greater total score are ranked ahead of those with lesser score.
Players with the same number of wins, losses and total score are
equivalently ranked or tied.

\head Example 1\endhead

Suppose that at the end of the penultimate round of a tournament,
the top four players (referred to as players $1$, $2$, $3$ and $4$
have records 
$(w,N_r-1-w,s)$,
$(w,N_r-1-w,s-d_1)$,
$(w,N_r-1-w,s-d_2)$,
and
$(w-2,N_r-w+2,s-d_3)$.
For brevity, we assume all players have played $N_r-1$ games 
and record only the differences in a matrix as follows:
$$
\pmatrix
0&0&0&2\cr
0&d_1&d_2&d_3\cr
\endpmatrix
$$

If we want to pair the top four players amongst themselves in the last
round, we have three choices of opponent for player $1$, each of which
then forces the remaining pairing. 
We want to compute the chance of each player being victorious
depending on the choice of pairing.

\subhead Player 1 vs.~player 2 \endsubhead

Player $1$ is the sole overall victor
(i) if in the final round player $1$ wins or ties and player $3$ loses;
(ii) if player $1$ wins, player $3$ wins, but player $1$'s final
total score is greater than player $3$'s;
or (iii) if player $1$ and player $3$ both tie.
The probability of this happening is therefore
$$
(1-D_{1,2}(-1))D_{3,4}(-1)
+\sum ?
+ P_{1,2}(0) P_{3,4}(0)
$$

\subhead Player 1 vs.~player 3 \endsubhead

\subhead Player 1 vs.~player 4 \endsubhead

% start from the premise that a suitable utility function for evaluating a pairing
 %system is one which measures the expected distribution of prize dollars and
 %%the players to whom they are awarded.  A simple example of such a function
 %would be the expected number of dollars that are awarded to the right player;
 %a more complex one would partially reward near misses in prize assignments.

% Spearman Rank Correlation Coefficient
\enddocument
